\documentclass[11pt]{article}
% ===========================================================================
%  Shared preamble for the AP Calculus AB exam-preparation worksheets.
%  Input right after \documentclass. Files beginning with "_" are not compiled
%  on their own. Navy and gold, to match the course's lesson pages.
% ===========================================================================
\usepackage[letterpaper,top=2.3cm,bottom=2.2cm,left=1.7cm,right=1.7cm,headheight=28pt,footskip=30pt]{geometry}
\usepackage{amsmath,amssymb}
\usepackage[T1]{fontenc}
\usepackage[utf8]{inputenc}
\usepackage{lmodern}
\usepackage{xcolor}
\usepackage[most]{tcolorbox}
\usepackage{enumitem}
\usepackage{fancyhdr}
\usepackage{array}
\usepackage{tikz}
\usepackage{pgfplots}
\pgfplotsset{compat=1.18}

\definecolor{navy}{HTML}{161E54}
\definecolor{ink}{HTML}{27357E}
\definecolor{gold}{HTML}{B98A00}
\definecolor{goldbg}{HTML}{FFFDF3}
\definecolor{inkbg}{HTML}{F6F7FD}
\definecolor{rule}{HTML}{C9CFF0}
\definecolor{okgreen}{HTML}{1B7A44}
\definecolor{okbg}{HTML}{F1F8F3}

\pagestyle{fancy}
\fancyhf{}
\fancyhead[L]{\footnotesize Name: \underline{\hspace{4cm}}}
\fancyhead[C]{\textbf{\color{navy}Integration Academy}\\[-2pt]{\scriptsize Calculus AB \textperiodcentered{} AP\textsuperscript{\textregistered} Exam Preparation}}
\fancyfoot[L]{\tiny\color{black!55}\parbox[t]{0.78\textwidth}{AP\textsuperscript{\textregistered} is a registered trademark of the College Board, which was not involved in the production of, and does not endorse, this material. \textcopyright{} Integration Academy.}}
\fancyfoot[R]{\scriptsize\thepage}
\renewcommand{\headrulewidth}{1.2pt}
\renewcommand{\headrule}{\hbox to\headwidth{\color{gold}\leaders\hrule height \headrulewidth\hfill}}
\renewcommand{\footrulewidth}{0pt}
\setlength{\parindent}{0pt}
\setlength{\parskip}{4pt}

% Title block: \wstitle{code}{title}{unit line}
\newcommand{\wstitle}[3]{\begin{center}{\small\bfseries\color{gold}WORKSHEET #1}\\[3pt]{\LARGE\bfseries\color{navy}#2}\\[3pt]{\small\color{black!60}#3}\end{center}\vspace{2pt}}
% Section heading
\newcommand{\sech}[1]{\par\vspace{10pt}{\large\bfseries\color{navy}#1}\par\nobreak\vspace{2pt}{\color{rule}\hrule height0.8pt}\vspace{7pt}}
% A level tag printed beside a group of questions
\newcommand{\level}[1]{\par\vspace{6pt}{\footnotesize\bfseries\color{gold}\MakeUppercase{#1}}\par\nobreak\vspace{2pt}}

\newtcolorbox{ideabox}{colframe=ink,colback=inkbg,coltitle=white,colbacktitle=ink,fonttitle=\bfseries\small,title={Key ideas},arc=2pt,boxrule=0.8pt,left=8pt,right=8pt}
\newtcolorbox{example}[1]{colframe=rule,colback=white,coltitle=navy,colbacktitle=inkbg,fonttitle=\bfseries\small,title={#1},arc=2pt,boxrule=0.8pt,borderline west={3pt}{0pt}{ink},left=10pt,right=8pt}
\newtcolorbox{frbox}[1]{colframe=gold,colback=goldbg,coltitle=white,colbacktitle=gold,fonttitle=\bfseries\small,title={#1},arc=2pt,boxrule=0.8pt,breakable,left=8pt,right=8pt}
\newtcolorbox{solbox}{colframe=okgreen,colback=okbg,coltitle=white,colbacktitle=okgreen,fonttitle=\bfseries\small,title={Worked solutions},arc=2pt,boxrule=0.8pt,breakable,left=8pt,right=8pt}
\newtcolorbox{scorebox}[1]{colframe=navy,colback=white,coltitle=white,colbacktitle=navy,fonttitle=\bfseries\small,title={#1},arc=2pt,boxrule=0.8pt,breakable,left=8pt,right=8pt}

\newcommand{\soln}{\par\smallskip\textbf{\color{ink}Solution.}\ }
\newcommand{\result}[1]{\par\smallskip\textbf{#1}}
% blank working space of a given height
\newcommand{\work}[1]{\par\vspace{#1}}
% a point value in the margin of a scoring guide
\newcommand{\pts}[1]{\hfill{\footnotesize\bfseries\color{gold}[#1]}}

% Coordinate grid for graph questions: \begin{apgraph}{xmin}{xmax}{ymin}{ymax} ... \end{apgraph}
\newenvironment{apgraph}[4]{\begin{center}\begin{tikzpicture}\begin{axis}[
  axis lines=middle, xmin=#1, xmax=#2, ymin=#3, ymax=#4,
  grid=both, grid style={black!12}, xtick distance=1, ytick distance=1,
  ticklabel style={font=\scriptsize}, xlabel={$x$}, ylabel={$y$},
  width=10cm, height=7.2cm, samples=80, clip=false,
  every axis plot/.append style={line width=1.1pt, color=ink}]}{\end{axis}\end{tikzpicture}\end{center}}
\newcommand{\closeddot}[2]{\addplot[only marks, mark=*, mark size=2.2pt, color=ink] coordinates {(#1,#2)};}
\newcommand{\opendot}[2]{\addplot[only marks, mark=*, mark size=2.2pt, color=ink, mark options={fill=white}] coordinates {(#1,#2)};}

% Multiple-choice options: four across, or stacked when the options are long.
\newcommand{\opts}[4]{\par\vspace{3pt}\begin{tabular}{@{}p{0.235\linewidth}p{0.235\linewidth}p{0.235\linewidth}p{0.235\linewidth}@{}}(A)\ #1 & (B)\ #2 & (C)\ #3 & (D)\ #4\end{tabular}\par\vspace{9pt}}
\newcommand{\optsv}[4]{\par\vspace{3pt}\begin{tabular}{@{}l@{\ \ }p{0.9\linewidth}@{}}(A) & #1\\(B) & #2\\(C) & #3\\(D) & #4\end{tabular}\par\vspace{9pt}}

\fancyhead[R]{\footnotesize\color{black!60}Worksheet 5.6}
\begin{document}

\wstitle{5.6}{Concavity and Points of Inflection}{Unit 5: Analysing Functions}

\begin{ideabox}
\begin{itemize}[leftmargin=*,itemsep=2pt,topsep=2pt]
  \item The graph of $f$ is \textbf{concave up} where $f''(x)>0$ (equivalently, where $f'$ is increasing) and \textbf{concave down} where $f''(x)<0$.
  \item A \textbf{point of inflection} is a point of the graph where the concavity changes.
  \item $f''(c)=0$ is not enough: $f''$ must \textbf{change sign} at $c$.
  \item A point of inflection is a point: give both coordinates when asked for the point.
\end{itemize}
\end{ideabox}

\sech{Worked examples}
\begin{example}{Example 1 \textperiodcentered{} A cubic}
Find the point of inflection of $f(x)=x^3+3x^2$.\soln $f''(x)=6x+6$ changes from negative to positive at $x=-1$.\result{The point $(-1,2)$.}
\end{example}

\begin{example}{Example 2 \textperiodcentered{} A quartic}
Find the $x$-coordinates of the points of inflection of $f(x)=x^4-12x^2$.\soln $f''(x)=12x^2-24=12\left(x^2-2\right)$ changes sign at $x=\pm\sqrt2$.\result{$x=-\sqrt2$ and $x=\sqrt2$.}
\end{example}

\begin{example}{Example 3 \textperiodcentered{} A bell-shaped curve}
Find the $x$-coordinates of the points of inflection of $f(x)=e^{-x^2}$.\soln $f'(x)=-2xe^{-x^2}$ and $f''(x)=\left(4x^2-2\right)e^{-x^2}$.\result{$x=\pm\dfrac{1}{\sqrt2}$.}
\end{example}

\newpage
\sech{Practice}

\level{Core}
\begin{enumerate}[leftmargin=*,itemsep=2pt]
  \item Show that the graph of $f(x)=x^2+3x$ is concave up everywhere.\work{2.3cm}
  \item Find the point of inflection of $f(x)=x^3$.\work{2.3cm}
  \item Describe the concavity of $f(x)=\ln x$.\work{2.3cm}
  \item A function has $f''(x)=x-1$. Where is its graph concave up?\work{2.3cm}
\end{enumerate}

\level{Exam level}
\begin{enumerate}[leftmargin=*,itemsep=2pt,start=5]
  \item Find the $x$-coordinates of the points of inflection of $f(x)=x^4-8x^3$.\work{2.8cm}
  \item Find the $x$-coordinate of the point of inflection of $f(x)=x^2\ln x$ for $x>0$.\work{2.8cm}
  \item A function has $f''(x)=(x+2)^2(x-3)$. Where does its graph have a point of inflection?\work{2.8cm}
  \item Find the point of inflection of $f(x)=\sin x$ on $(0,2\pi)$, and describe the concavity on each side.\work{2.8cm}
\end{enumerate}

\begin{frbox}{Question 9 \textperiodcentered{} Free-response practice}
Let $f(x)=x^5-10x^3$.
\begin{enumerate}[label=(\alph*),leftmargin=*,itemsep=2pt]
  \item Find $f''(x)$ in factored form.\work{2.6cm}
  \item Find the intervals on which the graph of $f$ is concave up.\work{2.6cm}
  \item Find the $x$-coordinates of all points of inflection. Justify your answer.\work{2.6cm}
\end{enumerate}
\end{frbox}

\newpage
\begin{solbox}
\begin{enumerate}[leftmargin=*,itemsep=4pt]
  \item $f''(x)=2>0$ for all $x$.
  \item $f''(x)=6x$ changes sign at $x=0$: the point $(0,0)$.
  \item $f''(x)=-\dfrac{1}{x^2}<0$: concave down for all $x>0$.
  \item For $x>1$.
  \item $f''(x)=12x^2-48x=12x(x-4)$ changes sign at $x=0$ and $x=4$.
  \item $f'(x)=2x\ln x+x$ and $f''(x)=2\ln x+3$, which changes sign at $x=e^{-3/2}$.
  \item Only at $x=3$: $(x+2)^2$ does not change sign.
  \item $f''(x)=-\sin x$ changes sign at $x=\pi$: the point $(\pi,0)$. Concave down on $(0,\pi)$, concave up on $(\pi,2\pi)$.
  \item \begin{enumerate}[label=(\alph*),leftmargin=*,itemsep=3pt]
    \item $f'(x)=5x^4-30x^2$ and $f''(x)=20x^3-60x=20x\left(x^2-3\right)$.
    \item $f''>0$ on $\left(-\sqrt3,0\right)$ and $\left(\sqrt3,\infty\right)$.
    \item $x=-\sqrt3$, $x=0$ and $x=\sqrt3$: $f''$ changes sign at each.
  \end{enumerate}
\end{enumerate}
\end{solbox}

\end{document}
